\chapter{Simulator Architecture}
\label{ch:architecture}

The prototype could afford to keep its state inside one agent because it represented one
vehicle and one task. That stopped working as soon as the right-side configuration panel was
introduced. A value could now exist in three forms: what the user was editing, what the
simulator had accepted, and what an already running episode was using. Resolving that ambiguity
became the central architectural problem. The final design follows the life of an experiment
from an editable idea to a frozen result.

\section{The simulator as an experimental product}

The application is organized around one user journey. The user configures a vehicle and task,
checks the derived policy interface, launches training or inference, and inspects the resulting
telemetry and artifacts. Figure~\ref{fig:product-workflow} shows this path. The final arrow back
to configuration is important: most useful changes in this project were responses to evidence
from the previous run.

\begin{figure}[htbp]
  \centering
  \begin{tikzpicture}[node distance=0.32cm]
    \node[rocketbox,text width=0.175\textwidth,minimum height=2.55cm] (configure) {
      \textcolor{rocketBlue}{\Large\textbf{1}}\\[2pt]
      \textbf{Configure}\\[3pt]
      \footnotesize vehicle, task, rewards, environment, trainer, faults and telemetry};
    \node[rocketbox,text width=0.175\textwidth,minimum height=2.55cm,right=of configure] (freeze) {
      \textcolor{rocketBlue}{\Large\textbf{2}}\\[2pt]
      \textbf{Check \& freeze}\\[3pt]
      \footnotesize validate feasibility, derive the policy schema and create an immutable session};
    \node[rocketbox,text width=0.175\textwidth,minimum height=2.55cm,right=of freeze] (execute) {
      \textcolor{rocketTeal}{\Large\textbf{3}}\\[2pt]
      \textbf{Execute}\\[3pt]
      \footnotesize train with PPO, run ONNX inference, evaluate or perform component tests};
    \node[rocketbox,text width=0.175\textwidth,minimum height=2.55cm,right=of execute] (inspect) {
      \textcolor{rocketAmber}{\Large\textbf{4}}\\[2pt]
      \textbf{Inspect}\\[3pt]
      \footnotesize review trajectories, reward terms, outcomes, checkpoints and manifests};
    \draw[rocketarrow] (configure) -- (freeze);
    \draw[rocketarrow] (freeze) -- (execute);
    \draw[rocketarrow] (execute) -- (inspect);
    \draw[-{Latex[length=2.3mm]},thick,dashed,draw=rocketAmber]
      ([yshift=-4mm]inspect.south) -- ++(0,-0.45) -| ([yshift=-4mm]configure.south)
      node[pos=0.47,below,font=\footnotesize,text=rocketInk]{evidence informs the next revision};
  \end{tikzpicture}
  \caption{The experiment workflow. Editable settings become a validated, immutable session
  before any simulation area or trainer consumes them.}
  \label{fig:product-workflow}
\end{figure}

The configuration panel mirrors the decisions needed for an experiment: vehicle hardware,
task and curriculum, rewards and termination, environment, learning, faults and telemetry.
Built-in presets provide a usable starting point, but their values remain editable and a user
can save new presets. The same path can therefore create a lightweight single-engine vehicle or
a Falcon~9-inspired multi-engine vehicle without placing vehicle-specific logic inside a task.

This workflow also separates four questions that the prototype used to blur together. Is the
configuration syntactically valid? Is the requested manoeuvre physically feasible? Does a
trained model still match the selected hardware? Does it actually succeed? Validation addresses
the first, component tests support the second, schema checks protect the third, and only a
frozen evaluation can answer the fourth.

\section{From editable settings to immutable state}

The first major redesign established a one-way configuration flow. The UI edits a
\texttt{SimulationSessionDraft}. Starting a run validates that draft and converts it into an
immutable \texttt{SimulationSessionSnapshot}. The \texttt{SimulationRunCoordinator} then creates
the requested simulation areas, and each \texttt{SimulationAreaHost} receives its own runtime
copy. From that point onward, the rocket and agent read the snapshot rather than the UI.

This boundary was added after scripts began reading and changing shared component values
directly. It was no longer clear whether a number on screen was a proposed value, a corrected
value after validation, or the value already used by training. Snapshotting turns ``start'' into
an explicit state transition. It also makes the accepted configuration serializable and gives
the manifest a stable object to hash.

The snapshot contains five main sections:

\begin{description}
  \item[Vehicle] geometry, mass, fuel, engine layout and actuator parameters;
  \item[Environment] task, atmosphere, wind, curricula, inference mode, seeds and faults;
  \item[Objective] task-specific rewards, shaping ranges and terminations;
  \item[Learning] behaviour schema, PPO settings and resume policy; and
  \item[Telemetry] enabled streams, sample rate and output location.
\end{description}

Validation derives dependent quantities and reports infeasible combinations before launch. A
warning can still be overridden for an intentional sandbox experiment, but that override is
stored in the session. It no longer disappears as an undocumented click.

\section{Breaking apart the original agent}

Once configuration had a clear owner, runtime behaviour needed the same treatment. The original
\texttt{FalconAgent} handled observations, actions, actuation, rewards, spawning, termination and
visual effects. It was convenient while the rocket was a cylinder; later, almost every change
touched the same file and created side effects elsewhere.

The current \texttt{RocketAssembly} builds the rigid body and selected components. Engines,
grid fins, reaction-control thrusters and landing legs own their local state, central systems
combine forces and mass properties, and the agent supplies normalized commands. Agent logic is
split by responsibility into lifecycle, observations, actuators, guidance, spawning, rewards,
telemetry and scenario-specific behaviour. Separate services own session lifecycle, area
creation and trainer startup. Some large partial classes remain, but new task behaviour now has
a defined location instead of another switch in a global script.

\section{A policy interface that follows the vehicle}

Modular hardware created the next problem: changing the vehicle also changes what the policy
must observe and control. Manually updating ML-Agents behaviour parameters became both tedious
and unsafe. The simulator now derives those dimensions from the accepted snapshot.

The base observation vector contains 21 values describing relative target kinematics,
orientation and global vehicle state. Engine timing adds five shared observations, followed by
eight values for each independently controlled engine: throttle, two gimbal coordinates, four
one-hot run states and a constraint timer. Leg landing adds four foot-contact values. Grid-fin,
RCS and separate engine-enable control add further component-dependent observations or actions.

For $N_e$ engine channels, $N_f$ fins and $N_r$ RCS channels, the continuous action size is
\begin{equation}
  n_a = (3+I_{\mathrm{enable}})N_e+N_f+N_r,
  \label{eq:actions}
\end{equation}
where $I_{\mathrm{enable}}\in\{0,1\}$. Policy outputs are bounded and then mapped to throttle,
gimbal or component-specific physical limits.

Shape alone is not enough to guarantee compatibility. Two observation vectors can have the same
length while assigning different meanings to their elements. The model artifact therefore
stores the derived schema, and resume or inference checks compare observation count, action
count, component layout, behaviour name, trainer type and network settings. This contract was
added after simulator changes made older policies unsafe to reuse and, in some cases, corrupted
the intended continuity of a run.

\section{Versioned runs instead of overwritten experiments}

The same problem appeared over time as well as across vehicles. A promising run often needed a
longer step budget, but resuming it by overwriting the old configuration erased the conditions
that had produced earlier checkpoints. The run format now keeps a top-level manifest and mutable
runtime progress alongside numbered, immutable revision directories. Every revision stores its
session JSON, launch manifest and trainer YAML; logs, checkpoints and the final ONNX model remain
associated with the run ID.

The launch manifest records the Python executable and package versions, GPU, seed, launch mode
and SHA-256 digest of the session. Curriculum counters are stored separately in runtime state
because they change while training proceeds. In this way a continuation can advance the same
policy without rewriting the identity of the experiment that created it. The L10 and L11
histories in Chapters~\ref{ch:development} and~\ref{ch:results} rely on this distinction.

\section{Starting Unity and the trainer in the right order}

The final architectural failure appeared before an episode had even begun. ML-Agents originally
opened its connection as soon as Unity entered Play mode, while the configurable rocket prefabs
had not yet been created. The trainer could therefore connect to an incomplete or unintended
agent.

The coordinator now validates the snapshot, creates every requested simulation area and only
then enables the ML-Agents behaviour. Dummy prefabs satisfy Unity editor references without
registering an agent; actual runtime prefabs are created after the configuration is known. The
same coordinator derives the final policy schema before allowing the run to start.

Physics advances every \SI{0.01}{s}. The standard decision period is three physics steps, so the
policy supplies a new command every \SI{0.03}{s}, while actuator state, forces and the contact
solver continue to update between decisions. This ordering keeps high-frequency physical
dynamics independent of the neural-network inference frequency and completes the path from an
editable experiment to a running, reproducible simulation.
